\chapter{Relative equations under positive lower mass}
\label{chap:relative-finite-measures}

For a finite non-pluripolar measure, the Monge--Amp\`ere equation can be solved in a prescribed relative envelope class. A positive exponent determines the potential uniquely; at exponent zero, the normalization scalar is unique, but the potential need not be. The proof passes from domination to compactness and then to approximation of the measure. Positive lower mass prevents the comparison envelopes from collapsing, while contact support controls their measures.

Fix a smooth semipositive big real $(1,1)$-form $\theta$ and a reference $\varphi$ satisfying \eqref{eq:als-model-condition}. Unless explicitly stated otherwise, assume
\begin{equation}\label{eq:relative-positive-lower-volume}
 v_-(\theta)>0.
\end{equation}
Fix $\varepsilon\in(0,1)$ with $P_{\varepsilon\theta}(\varphi)$ genuine, and write
\[
 \mathcal E_\varphi=\mathcal E_{\mathrm{env}}(X,\theta,\varphi),
 \qquad \mathcal E_0=\mathcal E_{\mathrm{env}}(X,\theta,0),
 \qquad \mu_u=\mathrm{MA}^{\mathrm{ext}}_\theta(u).
\]
Fix also a Hermitian form $\omega\geq\theta$ for capacity and a smooth positive volume form $\mathrm{d}V$ for $L^1$ convergence. The envelope constructions of \cref{sec:relative-envelope-room} remain available when members of $\mathcal E_\varphi$ have infinite mass. The solutions below have finite mass because their equations have finite right-hand sides.

\clearpage
\section{The general-measure equations}
\label{sec:relative-general-setting}

The exponential factor supplies comparison and fixes the additive constant. Letting the exponent tend to zero produces a second equation, with a scalar chosen by the geometry and the measure.

\Needspace{15\baselineskip}
\begin{theorem}
\label{thm:als-general-measure}
Let $\mu$ be a nonzero finite non-pluripolar measure. Then:
\begin{enumerate}
\item For every $\lambda>0$ there is a unique
$u\in\mathcal E_\varphi$ with $\mu_u=\mathrm{e}^{\lambda u}\mu$.
\item There are $u\in\mathcal E_\varphi$ and a unique $c>0$ with $\mu_u=c\mu$.
\end{enumerate}
\end{theorem}

The proof first treats the minimal reference in \cref{thm:general-measure-minimal}; those solutions then serve as comparison potentials in \cref{sec:general-measure-relative}. For general measure data, the conclusion is membership in $\mathcal E_\varphi$, rather than bounded difference from $\varphi$.

\section{Domination with one finite measure}
\label{sec:relative-finite-domination}

The first domination argument uses one finite measure. Contact support bounds the measures of stretched envelopes by that fixed measure, and its small tails exclude collapse. We use the stretching envelopes $Q_b$ of \cref{def:relative-envelope-class} and the uniform lower bound from \cref{thm:relative-uniform-mass-lower}:
\[
 \mu_a(X)\geq\kappa:=(1-\varepsilon)^n v_-(\theta)>0
 \qquad(a\in\mathcal E_\varphi).
\]

\begin{theorem}\label{thm:relative-finite-domination}
Let $u\in\mathcal E_\varphi$ have finite $\mu_u$, and let $v\in\PSH(X,\theta)$ satisfy $v\preceq\varphi$. If, for some $0\leq c<1$,
\[
 \mu_u\leq c\mu_v\quad\hbox{on }\{u<v\},
\]
then $u\geq v$.
\end{theorem}
\begin{proof}
Replace $v$ by $w=\max(u,v)$. Upward stability gives $w\in\mathcal E_\varphi$ and locality gives $\mu_u\leq c\mu_w$ on $\{u<w\}$. For $b>1$ put $q_b=Q_b(u,w)\in\mathcal E_\varphi$. It satisfies $q_b\leq u$, and its measure is carried by
\[
 D_b=\{b^{-1}q_b+(1-b^{-1})w=u\}.
\]
On $D_b$, positivity and ordered contact give
\[
 b^{-n}\mu_{q_b}+(1-b^{-1})^n\mu_w\leq\mu_u.
\]
Choose $b$ so large that $(1-b^{-1})^n>c$. On $D_b\cap\{u<w\}$ the preceding inequality forces both $\mu_w$ and $\mu_{q_b}$ to vanish. To justify this even when $\mu_w(X)=+\infty$, restrict first to common finite truncation levels of the potentials. On each such set the measures are finite and subtraction is legitimate. Exhausting the complements of the pole sets proves the assertion for the extended measures.

It follows that $\mu_{q_b}$ is carried by $D_b\cap\{u=w\}\subseteq\{q_b=u\}$. Since $q_b\leq u$, another application of ordered contact gives the bound with coefficient one
\begin{equation}\label{eq:relative-contact-fixed-measure}
 \mu_{q_b}\leq\mathbf1_{\{q_b=u\}}\mu_u.
\end{equation}
This coefficient is independent of $b$.

Along integers $b\to\infty$, the potentials $q_b$ decrease. If they collapsed identically to $-\infty$, \cref{thm:normalized-compactness} would give $M_b:=\sup_Xq_b\to-\infty$. But \eqref{eq:relative-contact-fixed-measure} would then imply
\[
 \kappa\leq\int_X\mu_{q_b}
 \leq\mu_u(\{u\leq M_b\})\longrightarrow0.
\]
Here finiteness of $\mu_u$ and its vanishing on the pole set are exactly what gives the limit. This is impossible. Hence the decreasing limit $q$ is genuine. On $\{u<w-a\}$, for any $a>0$, the defining inequality gives $q_b\leq u-(b-1)a$. Outside the poles its limit is $-\infty$ there. Such a set has smooth volume zero because $q$ is genuine. Thus $u\geq w$ almost everywhere, and then everywhere by \cref{lem:qpsh-almost-everywhere-order}. This proves the theorem.
\end{proof}

\subsection{Domination for extended measures}

To prove zero domination for possibly infinite measures, we construct a difference envelope with no mass where it is negative. This makes its entire measure finite and permits the preceding theorem to be applied. The resulting comparison yields domination for the original potentials.

\Needspace{10\baselineskip}
\begin{lemma}\label{lem:relative-zero-domination}
Let $u,v\in\mathcal E_\varphi$ satisfy $u\leq v$. If
\[
 \mu_u(\{u<v\})=0,
\]
then $u=v$.
\end{lemma}
\begin{proof}
Translate both potentials by a common constant so that $u\leq v\leq\varphi\leq0$. Fix $\varepsilon<\delta<1$, put $\beta=1-\delta$, and set
\[
 q=P_{\delta\theta}(v).
\]
This envelope is genuine by \cref{thm:relative-uniform-mass-lower}. By \cref{lem:relative-difference-envelope},
\[
 w=P_\theta(u-\varphi)\in\mathcal E_0,
 \qquad w+\varphi\leq u,\qquad w\leq0.
\]
The stretching lemma with reference zero shows that
$z=P_\theta(w/\beta)$ belongs to $\mathcal E_0$. Scaling the background and each defining competitor gives
\[
 \beta z\in\mathcal E_{\mathrm{env}}(X,\beta\theta,0).
\]
Indeed, if $a\in\PSH(X,\theta)$ lies below $t z$, then $\beta a\in\PSH(X,\beta\theta)$ lies below $t\beta z$, for every $t\geq1$.

Since $q\leq v\leq\varphi$, we have
\[
 \beta z\leq w\leq u-\varphi\leq u-q
\]
off the pole sets. Thus
\[
 k=P_{\beta\theta}(u-q)
\]
is genuine. Here the difference obstacle is taken off the pluripolar pole sets; it is quasi-continuous and finite almost everywhere. The general envelope theorem gives a finite upper bound for $k$ and the canonical inequality $k+q\leq u$. Upward stability therefore gives
$k\in\mathcal E_{\mathrm{env}}(X,\beta\theta,0)$.

Write $\nu_k=\mathrm{MA}^{\mathrm{ext}}_{\beta\theta}(k)$. Contact support carries $\nu_k$ by
\[
 D=\{k+q=u\}
\]
outside the pole sets. Positivity of mixed products and ordered contact give
\[
 \nu_k\leq\mathbf1_D\mu_u.
\]
These inequalities are first restricted to common finite truncation levels and then exhausted. On $D\cap\{k<0\}$ we have $u=k+q<q\leq v$, so $\nu_k(\{k<0\})=0$.

This vanishing proves that $\nu_k$ is finite. In fact, $\widehat k=\max(k,-1)$ is a single bounded $\beta\theta$-plurisubharmonic function, and locality gives
\[
 \nu_k(X)=\nu_k(\{k\geq0\})
 =\int_{\{k\geq0\}}\mathrm{MA}_{\beta\theta}(\widehat k)
 \leq\int_X\mathrm{MA}_{\beta\theta}(\widehat k)<\infty.
\]

Apply \cref{thm:relative-finite-domination} in the background $\beta\theta$, with reference zero and potentials $k,0$, using coefficient zero. Its hypotheses hold: $\beta\theta$ is semipositive and big, zero satisfies the model condition, and
\[
 v_-(\beta\theta)=\beta^n v_-(\theta)>0.
\]
The scaling identity follows by the bijection $a\mapsto\beta a$ between bounded potentials in the two backgrounds. We have just proved finiteness of $\nu_k$. The theorem yields $k\geq0$, hence $u\geq q$.

To pass to $\delta\uparrow1$, fix $\delta_0\in(\varepsilon,1)$ and $r=P_{\delta_0\theta}(v)$. For $\delta_0<\delta<1$, set $t=(\delta-\delta_0)/(1-\delta_0)$. The potential $t v+(1-t)r$ has background $\delta\theta$ and lies below $v$. Therefore
\[
 u\geq P_{\delta\theta}(v)\geq t v+(1-t)r.
\]
Let $t\uparrow1$ at points where $v,r$ are finite. This gives $u\geq v$ almost everywhere, and then everywhere by \cref{lem:qpsh-almost-everywhere-order}. Together with $u\leq v$, it proves equality.
\end{proof}

\begin{theorem}\label{thm:relative-extended-domination}\index{Domination principle!extended measures}
Let $u\in\mathcal E_\varphi$ and $v\in\PSH(X,\theta)$ with $v\preceq\varphi$. If $0\leq c<1$ and
\[
 \mu_u\leq c\mu_v\quad\hbox{on }\{u<v\},
\]
then $u\geq v$.
\end{theorem}
\begin{proof}
Put $w=\max(u,v)\in\mathcal E_\varphi$ and choose $b>1$ with $(1-b^{-1})^n>c$. The stretching lemma gives $p=Q_b(u,w)\in\mathcal E_\varphi$ and $p\leq u$. Its measure is carried by
\[
 D=\{b^{-1}p+(1-b^{-1})w=u\}.
\]
On $D$, positivity and ordered contact imply
\[
 b^{-n}\mu_p+(1-b^{-1})^n\mu_w\leq\mu_u.
\]
On $\{u<w\}$, locality gives $\mu_u\leq c\mu_w$. Restricting both inequalities to $D\cap\{u<w\}$ and to common finite truncation levels makes subtraction legitimate. Since $(1-b^{-1})^n>c$, it follows that $\mu_p$ vanishes there. Exhausting the finite levels proves that $\mu_p$ is carried by $D\cap\{u=w\}\subseteq\{p=u\}$. Thus $\mu_p(\{p<u\})=0$.

Apply \cref{lem:relative-zero-domination} to $p\leq u$, both in $\mathcal E_\varphi$. It yields $p=u$. Substitution into the defining convex inequality for $p$ gives $w\leq u$ off the poles, hence everywhere by \cref{lem:qpsh-almost-everywhere-order}. Thus $v\leq w\leq u$.
\end{proof}

\begin{corollary}\label{cor:relative-extended-exponential-order}\label{cor:relative-finite-exponential-order}
Under the hypotheses on $u,v$ in \cref{thm:relative-extended-domination}, for every $\lambda>0$,
\[
 \mathrm{e}^{-\lambda u}\mu_u\leq \mathrm{e}^{-\lambda v}\mu_v
 \quad\Longrightarrow\quad u\geq v.
\]
Also $\mu_u\leq c\mu_v$ on $X$ for a finite constant $c\geq0$ implies $c\geq1$.
\end{corollary}
\begin{proof}
On $\{u<v-a\}$ the exponential inequality gives $\mu_u\leq \mathrm{e}^{-\lambda a}\mu_v$. Apply the theorem to $v-a$ and let $a\downarrow0$. All weights are taken off the pluripolar pole sets. For the second assertion, if $c<1$, applying the theorem to $v+A$ for every real $A$ contradicts finiteness of both potentials at a common point.
\end{proof}

Applying the exponential implication in both directions proves uniqueness for a fixed positive exponent. Applied to two equations $\mu_u=c_1\nu$ and $\mu_v=c_2\nu$, the scalar assertion gives $c_1=c_2$. The final section shows why this does not determine the normalized potential at exponent zero.


\section{Common envelopes and dominated continuity}

\label{sec:relative-common-envelope-setting}
Let $u_j\in\mathcal E_\varphi$ converge in $L^1$, with measures dominated by one finite non-pluripolar measure $\nu$. To keep the limit in $\mathcal E_\varphi$, we construct, after passing to a subsequence, a common lower envelope in this class. Partitioning the contact set of a finite rooftop according to the obstacle it touches shows that the rooftop measure is again dominated by $\nu$, with no factor depending on the number of obstacles. A summable integration error, combined with the lower mass bound $\kappa$, then prevents the stretched rooftops from collapsing, uniformly in the number of obstacles.

\begin{lemma}\label{lem:dominated-relative-rooftop-extraction}\index{Envelope!common lower envelope}
Suppose $u_j\in\mathcal E_\varphi$, $u_j\to u\in\PSH(X,\theta)$ in $L^1(X)$, and $\mu_{u_j}\leq\nu$ for a fixed finite non-pluripolar measure $\nu$. There is a subsequence, denoted again by $u_j$, such that
\[
 h=P_\theta\left(\inf_j u_j\right)\in\mathcal E_\varphi.
\]
In particular $u\in\mathcal E_\varphi$.
\end{lemma}
\begin{proof}
We first bound stretched rooftops of the tail of the sequence from below and then add the finitely many remaining terms.

The suprema of an $L^1$-convergent sequence are bounded, so \eqref{eq:relative-normalized-model-bound} gives a constant $C$ with $u_j\leq\varphi+C$ for all $j$. After a common translation we may assume $u_j\leq\varphi\leq0$, and then $u\leq\varphi$. Since $\mu_{u_j}\leq\nu$, \cref{lem:fixed-npp-integrals} gives
\[
 I_j=\int_X|\mathrm{e}^{u_j}-\mathrm{e}^u|\,\mathrm{d}\mu_{u_j}\longrightarrow0.
\]
We extract, once and for all, a subsequence with $I_j\leq2^{-j}$ along which $u_j\to u$ almost everywhere for smooth volume.

\Needspace{8\baselineskip}
For $N\geq j$ put $p_{j,N}=P_\theta(u_j,\ldots,u_N)$. These rooftops belong to $\mathcal E_\varphi$ by \cref{prop:relative-finite-rooftops}. Partition their contact sets according to the first index at which the minimum is attained. Ordered contact on these disjoint pieces gives both
\begin{equation}\label{eq:finite-rooftop-common-domination}
 \mu_{p_{j,N}}\leq\nu,
 \qquad
 \int_X|\mathrm{e}^{p_{j,N}}-\mathrm{e}^u|\,\mathrm{d}\mu_{p_{j,N}}
 \leq\sum_{\ell=j}^N I_\ell\leq2^{1-j}.
\end{equation}
The first bound has no factor $N$, because each piece is controlled by the restriction of the same measure $\nu$ to that piece.

Fix a real number $B\geq1$ and set $r_{j,N}=Q_B(p_{j,N},\varphi)$. It belongs to $\mathcal E_\varphi$ by the stretching lemma, so $\mu_{r_{j,N}}(X)\geq\kappa$. With $D_{j,N}=\{B^{-1}r_{j,N}+(1-B^{-1})\varphi=p_{j,N}\}$, contact support and positivity give
\[
 \mu_{r_{j,N}}\leq B^n\mathbf1_{D_{j,N}}\mu_{p_{j,N}}
 \leq B^n\nu.
\]
On this contact set, $\varphi\leq0$ gives
$\mathrm{e}^{p_{j,N}}\leq \mathrm{e}^{r_{j,N}/B}$. Consequently
\begin{equation}\label{eq:relative-rooftop-noncollapse-test}
 \int_X\mathrm{e}^u\,\mathrm{d}\mu_{r_{j,N}}
 \leq B^n2^{1-j}
      +B^n\nu(X)\exp\!\left(\frac{\sup_X r_{j,N}}B\right).
\end{equation}
On the other hand, choose $K>0$ with $B^n\nu(\{u\leq-K\})<\kappa/2$; this is possible because $u\not\equiv-\infty$ and $\nu$ does not charge its pole set. The measure domination then gives $\mu_{r_{j,N}}(\{u\leq-K\})<\kappa/2$, and the lower mass bound yields
\[
 \int_X\mathrm{e}^u\,\mathrm{d}\mu_{r_{j,N}}\geq \mathrm{e}^{-K}\kappa/2.
\]
Choose $j$ so large that $B^n2^{1-j}<\mathrm{e}^{-K}\kappa/4$. Then \eqref{eq:relative-rooftop-noncollapse-test} bounds $\sup_Xr_{j,N}$ from below, independently of $N$. Since the sequence decreases in $N$ and is bounded above by $\varphi$, its limit $R$ is not identically $-\infty$. It satisfies, for every $\ell\geq j$,
\begin{equation}\label{eq:relative-tail-stretched-competitor}
 B^{-1}R+(1-B^{-1})\varphi\leq u_\ell,
 \qquad R\leq\varphi.
\end{equation}
The second inequality follows already from $r_{j,N}\leq p_{j,N}\leq\varphi$.

It remains to include the finitely many earlier terms and to verify every defining condition of the relative class. Fix $b\geq1$, apply the preceding construction with $B=2b$, and let $j$ and $R$ be as above. If $j>1$, let $p_0=P_\theta(u_1,\ldots,u_{j-1})$ and $S=Q_{2b}(p_0,\varphi)$. By the finite-rooftop and stretching lemmas, neither is identically $-\infty$, and $S\leq\varphi$. If $j=1$, set $S=\varphi$. The potential $H=(R+S)/2$ is $\theta$-plurisubharmonic. For a tail index $\ell\geq j$ we use \eqref{eq:relative-tail-stretched-competitor} and $S\leq\varphi$; for an earlier index we use the analogous inequality for $S$ and $R\leq\varphi$. In either case,
\[
 b^{-1}H+(1-b^{-1})\varphi\leq u_\ell
 \quad\hbox{for every }\ell\geq1.
\]
For $b=1$ this shows that $h=P_\theta(\inf_\ell u_\ell)$ is not identically $-\infty$. For general $b$ the left side is a competitor for $h$, so the same inequality holds with $h$ in place of every $u_\ell$; this is the defining condition of the relative class for $b$. Also $h\leq\varphi$. Thus $h\in\mathcal E_\varphi$.

Finally $h\leq u_j$ for all $j$, so the almost-everywhere convergence and \cref{lem:qpsh-almost-everywhere-order} give $h\leq u\leq\varphi$. Upward stability proves $u\in\mathcal E_\varphi$. Undoing the common translation proves the lemma as stated.
\end{proof}

Once a common lower envelope is available, convergence in capacity follows from a single integral condition. The proof uses the zero-domination lemma, so infinite masses are allowed.

\Needspace{17\baselineskip}
\begin{proposition}\label{prop:relative-common-envelope-criterion}
Let $h\in\mathcal E_\varphi$, let $C$ be a constant, and let $u_j\in\PSH(X,\theta)$ satisfy
\[
 h\leq u_j\leq\varphi+C,\qquad
 u_j\longrightarrow u\in\PSH(X,\theta)\quad\hbox{in }L^1(X).
\]
If
\[
 I_j:=\int_X|\mathrm{e}^{u_j}-\mathrm{e}^u|\,\mu_{u_j}\longrightarrow0,
\]
then $u\in\mathcal E_\varphi$ and $u_j\to u$ in capacity.
\end{proposition}
\begin{proof}
The bounds pass to $u$ almost everywhere along a subsequence, and then everywhere by \cref{lem:qpsh-almost-everywhere-order}; upward stability gives $u_j,u\in\mathcal E_\varphi$. After translating the entire family and its bounds by a common constant, we may assume that all potentials are nonpositive.

It suffices to prove convergence in capacity along a further subsequence of every subsequence. Extract so that $I_j\leq2^{-j}$ and $u_j\to u$ almost everywhere for smooth volume, and form
\[
 p_{j,N}=P_\theta(u_j,\ldots,u_N),\qquad
 p_j=P_\theta\left(\inf_{\ell\geq j}u_\ell\right).
\]
All these envelopes lie above $h$ and below $\varphi+C$; hence none of them is identically $-\infty$, and all belong to $\mathcal E_\varphi$. Partition the contact set of $p_{j,N}$ according to the first index $\ell$ with $p_{j,N}=u_\ell$. Ordered contact on each part gives
\[
 \int_X|\mathrm{e}^{p_{j,N}}-\mathrm{e}^u|\,\mu_{p_{j,N}}
 \leq\sum_{\ell=j}^N I_\ell\leq2^{1-j}.
\]
Only integrals of nonnegative functions occur here, so the estimate holds even if a rooftop has infinite mass.

As $N\to\infty$, $p_{j,N}\downarrow p_j$; as $j\to\infty$, the $p_j$ increase, and their regularized limit $p$ is not identically $-\infty$. Both monotone limits hold in capacity by \cref{lem:monotone-psh-capacity}. Moreover $h\leq p\leq u\leq\varphi+C$, where $p\leq u$ follows from the almost-everywhere convergence of the extracted $u_j$. Thus $p\in\mathcal E_\varphi$.

Apply \cref{lem:weighted-extended-lsc}, first as $N\to\infty$ and then as $j\to\infty$, with the bounded quasi-continuous weights in the preceding estimate. These weights converge in capacity because the exponential is Lipschitz on a common upper half-line. We obtain
\[
 \int_X|\mathrm{e}^p-\mathrm{e}^u|\,\mu_p=0,
 \qquad \mu_p(\{p<u\})=0.
\]
The second equality follows off the poles, which the extended measures ignore. The zero-domination lemma \cref{lem:relative-zero-domination} now gives $p=u$; no finiteness of $\mu_p$ is needed.

Finally, the upper envelopes $U_j=(\sup_{\ell\geq j}u_\ell)^*$ decrease to $u$, by Hartogs' lemma (\cref{lem:qpsh-hartogs}). Applying \cref{lem:monotone-psh-capacity} to $p_j\uparrow u$ and $U_j\downarrow u$, the inequalities $p_j\leq u_j\leq U_j$ give convergence in capacity along this subsequence, and hence for the original sequence.
\end{proof}

\subsection{Convergence under a fixed measure}

The extraction lemma provides the lower envelope required by the integral criterion. The same fixed measure then controls the measures of deep sublevel sets, so convergence in capacity passes to the Monge--Amp\`ere measures themselves.

\begin{theorem}\label{thm:relative-fixed-measure-continuity}\label{thm:relative-dominated-continuity}\index{Non-pluripolar measure!dominated continuity}
Suppose $u_j\in\mathcal E_\varphi$, $u_j\to u\in\PSH(X,\theta)$ in $L^1(X)$, and $\mu_{u_j}\leq\nu$ for a fixed finite non-pluripolar measure $\nu$. Then
\[
 u\in\mathcal E_\varphi,
 \qquad u_j\longrightarrow u\text{ in capacity},
 \qquad \mu_{u_j}\rightharpoonup\mu_u\leq\nu.
\]
\end{theorem}
\begin{proof}
Start with an arbitrary subsequence. By \cref{lem:dominated-relative-rooftop-extraction}, a further subsequence has a common lower envelope $h\in\mathcal E_\varphi$, and $u\in\mathcal E_\varphi$. Inequality \eqref{eq:relative-normalized-model-bound} and Hartogs' lemma give a common upper bound $u_j\leq\varphi+C$.

After a common translation, assume $u_j,u\leq0$. Since $\mu_{u_j}\leq\nu$, \cref{lem:fixed-npp-integrals} gives
\[
 \int_X|\mathrm{e}^{u_j}-\mathrm{e}^u|\,\mathrm{d}\mu_{u_j}
 \leq\int_X|\mathrm{e}^{u_j}-\mathrm{e}^u|\,\mathrm{d}\nu
 \longrightarrow0.
\]
Thus \cref{prop:relative-common-envelope-criterion} gives convergence in capacity along this subsequence. For the measures themselves, the common lower bound gives
\[
 \mu_{u_j}(\{u_j\leq-T\})
 \leq\nu(\{h\leq-T\})\longrightarrow0
 \quad\hbox{uniformly in }j,
\]
since $\nu$ is finite and does not charge the pole set of $h\not\equiv-\infty$. The tail criterion for smooth backgrounds, \cref{thm:general-background-damping}, therefore gives $\mu_{u_j}\rightharpoonup\mu_u$, and weak limits preserve domination by the fixed finite measure $\nu$. Every subsequence of the original sequence has a further subsequence with these conclusions and the same limit $u$, so both convergences hold for the original sequence.
\end{proof}


\section{Moving bounded carrier measures}

In the approximation scheme of \cref{sec:general-measure-relative}, the measures of the approximating solutions are dominated by Monge--Amp\`ere measures of uniformly bounded $\omega$-plurisubharmonic potentials $\psi_j$, which vary with $j$. For a bounded $\omega$-plurisubharmonic potential $\psi$, write $\omega_\psi=\omega+\ddc\psi$. A single dominating measure is then unavailable, but an exponential domination $\mu_{u_j}\leq A\mathrm{e}^{\lambda u_j}\omega_{\psi_j}^n$ still provides a common lower envelope. A smaller-background potential and a strict big potential give comparison functions whose constants are uniform in $j$. If every $u_j$ has the ordinary singularity type of $\varphi$, smaller-background comparison gives the same construction even when $v_-(\theta)=0$. Convergence of the potentials $\psi_j$ in capacity then identifies the limiting equation.

\Needspace{19\baselineskip}
\begin{proposition}\label{prop:exponential-common-rooftop}
Fix $A,\lambda>0$, $C\in\mathbb R$ and $L\geq0$. For $j\geq1$, suppose that $u_j\in\mathcal E_\varphi$, $u_j\leq\varphi+C$, and
\begin{equation}\label{eq:exponential-bounded-reference-domination}
 \mu_{u_j}\leq A \mathrm{e}^{\lambda u_j}\omega_{\psi_j}^n,
 \qquad \psi_j\in\PSH(X,\omega),\qquad
 \|\psi_j\|_\infty\leq L.
\end{equation}
Assume either $v_-(\theta)>0$ or $u_j\sim\varphi$ for every $j$. Then
\[
 h=P_\theta\left(\inf_j u_j\right)\in\mathcal E_\varphi.
\]
\end{proposition}

No convergence of $\psi_j$ or $u_j$ is assumed, and in the second alternative the constants in $u_j\sim\varphi$ need not be uniform in $j$.

\begin{proof}
Translate all $u_j$ by one constant so that $u_j\leq\varphi\leq0$; this changes $A$ by a fixed factor. Translate each $\psi_j$ so that it takes values in $[-2L,0]$; this does not change its measure. Set $q=P_{\varepsilon\theta}(\varphi)$ and choose $\rho\in\PSH(X,\theta)$ and $a>0$ such that
\[
 q\leq\varphi,
 \qquad \rho\leq0,\qquad \theta+\ddc\rho\geq a\omega.
\]
Here $q\not\equiv-\infty$ by the model condition, and $\rho$ exists because $\theta$ is big. Fix $0<\delta<(1-\varepsilon)a$.

\Needspace{8\baselineskip}
For a fixed real $b\geq1$, put $r_{j,b}=Q_b(u_j,\varphi)$. The stretching lemma puts it in $\mathcal E_\varphi$, and $r_{j,b}\leq u_j\leq0$. Contact support and positivity imply
\[
 \mu_{r_{j,b}}
 \leq b^n\mathbf1_{D_{j,b}}\mu_{u_j},
 \qquad
 D_{j,b}=\{b^{-1}r_{j,b}+(1-b^{-1})\varphi=u_j\}.
\]
On this contact set, $\varphi\leq0$ gives $u_j\leq r_{j,b}/b$. Thus \eqref{eq:exponential-bounded-reference-domination} yields
\begin{equation}\label{eq:stretched-exponential-domination}
 \mathrm{e}^{-\lambda r_{j,b}/b}\mu_{r_{j,b}}
 \leq b^n A\omega_{\psi_j}^n,
\end{equation}
where the weights are evaluated off the pluripolar pole sets.

Choose $C_b\geq0$ so large that
\[
 \delta^n\exp(\lambda C_b/b)\geq b^n A,
\]
and define
\[
 s_{j,b}=q+(1-\varepsilon)\rho+\delta\psi_j-C_b.
\]
This is a $\theta$-plurisubharmonic function, not identically $-\infty$, since
\[
 \begin{split}
 \theta+\ddc s_{j,b}
 &= (\varepsilon\theta+\ddc q)
      +(1-\varepsilon)(\theta+\ddc\rho)
      +\delta\ddc\psi_j\\
 &\geq \delta\omega_{\psi_j}.
 \end{split}
\]
Consequently $\mu_{s_{j,b}}\geq\delta^n\omega_{\psi_j}^n$: this inequality holds away from the poles by the local positive product inequality, and then globally because the bounded reference measure does not charge pluripolar sets. Also $s_{j,b}\leq-C_b$ and $s_{j,b}\preceq\varphi$. Together with \eqref{eq:stretched-exponential-domination}, we conclude
\[
 \mathrm{e}^{-\lambda s_{j,b}/b}\mu_{s_{j,b}}
 \geq b^n A\omega_{\psi_j}^n
 \geq \mathrm{e}^{-\lambda r_{j,b}/b}\mu_{r_{j,b}}.
\]

If $v_-(\theta)>0$, apply \cref{cor:relative-extended-exponential-order} with first potential $r_{j,b}$ and second potential $s_{j,b}$. It gives $r_{j,b}\geq s_{j,b}$.

In the second alternative, $r_{j,b}\sim\varphi$: if $u_j\geq\varphi-K_j$, the potential $\varphi-bK_j$ is a competitor for $r_{j,b}$. Hence the smaller-background envelope $P_{\varepsilon\theta}(r_{j,b})$ is not identically $-\infty$, and $s_{j,b}\preceq r_{j,b}$. Therefore \cref{cor:singular-exponential-order} also gives $r_{j,b}\geq s_{j,b}$, without a lower-mass hypothesis.

Since $\psi_j\geq-2L$, we now have a lower bound independent of $j$:
\[
 H_b=q+(1-\varepsilon)\rho-2\delta L-C_b
 \leq r_{j,b}.
\]
The potential $H_b$ is $\theta$-plurisubharmonic and not identically $-\infty$. By the definition of $r_{j,b}$, therefore,
\[
 b^{-1}H_b+(1-b^{-1})\varphi\leq u_j
 \quad\hbox{for every }j.
\]
For $b=1$ this shows that the infimum envelope $h$ is not identically $-\infty$. For general $b$, the left side is a competitor for $h$, so it lies below $h$; it is thus a competitor in the defining condition of the relative class for $b$. Since $h\leq\varphi$, we obtain $h\in\mathcal E_\varphi$. Undoing the common translation finishes the proof.
\end{proof}

\Needspace{8\baselineskip}
\begin{corollary}\label{cor:exponential-relative-class-closure}
Under the hypotheses of \cref{prop:exponential-common-rooftop}, if $u_j\to u\in\PSH(X,\theta)$ in $L^1(X)$, then $u\in\mathcal E_\varphi$.
\end{corollary}
\begin{proof}
The common envelope gives $h\leq u_j\leq\varphi+C$. Along an almost-everywhere convergent subsequence, \cref{lem:qpsh-almost-everywhere-order} gives $h\leq u\leq\varphi+C$, and upward stability gives the conclusion.
\end{proof}


\subsection{Convergence with a common envelope}

If the potentials $\psi_j$ converge in capacity, the integration error in \cref{prop:relative-common-envelope-criterion} tends to zero, and the common lower envelope controls the measures of deep sublevel sets. The dominating measure may also contain a fixed finite non-pluripolar part.

\Needspace{23\baselineskip}
\begin{theorem}\label{thm:relative-moving-measure-continuity}\index{Monge--Ampere measure@Monge--Amp\`ere measure!moving bounded carriers}
Let $h\in\mathcal E_\varphi$, let $C$ be a constant, and let $u_j\in\PSH(X,\theta)$ satisfy
\[
 h\leq u_j\leq\varphi+C,\qquad
 u_j\longrightarrow u\in\PSH(X,\theta)\quad\hbox{in }L^1(X).
\]
Fix $A\geq0$ and a finite non-pluripolar measure $\nu$ on $X$. Let $\psi_j\in\PSH(X,\omega)$ be uniformly bounded and converge in capacity to a bounded $\psi\in\PSH(X,\omega)$. Suppose
\[
 \mu_{u_j}\leq A\omega_{\psi_j}^n+\nu.
\]
Then $u\in\mathcal E_\varphi$, $u_j\to u$ in capacity, $\mu_u$ is finite, and
\[
 \mu_{u_j}\rightharpoonup\mu_u.
\]
\end{theorem}
\begin{proof}
Translate so that $u_j,u\leq0$. The functions $\mathrm{e}^{u_j},\mathrm{e}^u$ are uniformly bounded quasi-plurisubharmonic functions with a common lower curvature bound. They converge in $L^1$ for smooth volume, since the exponential is Lipschitz on $(-\infty,0]$. Hence \cref{lem:moving-bounded-integrals,lem:fixed-npp-integrals} give
\[
 \int_X|\mathrm{e}^{u_j}-\mathrm{e}^u|\,\omega_{\psi_j}^n\longrightarrow0,
 \qquad
 \int_X|\mathrm{e}^{u_j}-\mathrm{e}^u|\,\mathrm{d}\nu\longrightarrow0.
\]
Thus the integration error in \cref{prop:relative-common-envelope-criterion} tends to zero, and that proposition gives convergence in capacity and $u\in\mathcal E_\varphi$.

By \cref{lem:bounded-mixed-capacity}, the bounded reference measures $\omega_{\psi_j}^n$ have uniformly bounded mass and satisfy a common capacity estimate. Hence, for a constant $B$ independent of $j$ and $T$,
\[
 \mu_{u_j}(\{u_j<-T\})
 \leq B\operatorname{Cap}_\omega(\{h<-T\})
       +\nu(\{h<-T\})\longrightarrow0
\]
uniformly in $j$. The first term tends to zero because $h$ is a quasi-plurisubharmonic function not identically $-\infty$, and the second because $\nu$ is finite and does not charge the pole set of $h$. The total masses $\mu_{u_j}(X)$ are uniformly bounded by the assumed domination. The tail criterion for smooth backgrounds, \cref{thm:general-background-damping}, therefore gives finiteness of $\mu_u$ and the asserted weak convergence.
\end{proof}


\section{General measures in the minimal envelope class}
\label{sec:general-measure-minimal}

The bounded subsolution theorem and dominated continuity already give the equation for every finite non-pluripolar measure in the minimal envelope class. The missing link is a bounded reference adapted to the measure and to the same background. Its density need only be integrable. We construct that reference from a standard local decomposition theorem, and then give the limiting argument. This reconstructs the minimal-reference case of \cite[Theorem 5.8 and Corollary 5.9]{ALS26}.

\begin{lemma}\label{lem:bounded-measure-carrier}
Let $\theta\geq0$ be smooth and big, and let $\mu$ be a finite non-pluripolar measure on $X$. There are a bounded $\psi\in\PSH(X,\theta)$ and a nonnegative
$f\in L^1(\mathrm{MA}_\theta(\psi))$ such that
\[
 \mu=f\,\mathrm{MA}_\theta(\psi).
\]
No upper or lower volume property is assumed.
\end{lemma}
\begin{proof}
We first specify the local input. Cegrell's decomposition theorem \cite[Theorem 6.3]{Ceg98} says that a positive compactly supported measure $\tau$ on a bounded hyperconvex domain in $\mathbb C^n$, $n\geq2$, is
\[
 \tau=g(\ddc v)^n+\tau_s,
 \qquad v\in\mathcal E_0,\qquad
 0\leq g\in L^1((\ddc v)^n),
\]
where $\tau_s$ is carried by a pluripolar set. Here the local class $\mathcal E_0$ consists of bounded negative plurisubharmonic functions with zero boundary values and finite total Monge--Amp\`ere mass. This local notation is different from the global envelope class $\mathcal E_0$ in this chapter. If $\tau$ is non-pluripolar, $\tau_s=0$, so $\tau\ll(\ddc v)^n$. A change in the normalization of $\ddc$ only changes $g$ by a positive constant.

For completeness the required local assertion on a disk also follows from the Green potential. Let $a\leq0$ be the Green potential of a finite compactly supported non-polar measure $\tau$, normalized by $\ddc a=\tau$. Then $v=e^a$ is bounded and subharmonic, and
\[
 \ddc v\geq e^a\tau.
\]
To justify the inequality, set $a_M=\max(a,-M)$. The bounded subharmonic chain rule gives
$\ddc e^{a_M}\geq e^{a_M}\ddc a_M$, which on $\{a>-M\}$ is at least $e^a\tau$. Let $M\to\infty$, using decreasing convergence of the bounded subharmonic functions $e^{a_M}$ and monotone convergence on the right. Since $\{a=-\infty\}$ is polar, $e^a>0$ almost everywhere for $\tau$. Thus $\tau\ll\ddc v$. The bounded chain inequality itself follows by local convolution and passage to the distributional limit. We use only this absolute continuity assertion in every dimension.

Choose a strict potential $\rho$ with analytic singularities, smooth off its analytic pole set $Z$, and
\[
 \theta+\ddc\rho\geq\kappa\omega,\qquad\kappa>0.
\]
Such a potential is available by strict regularization. Cover $X\setminus Z$ by countably many compact inner coordinate balls $K_i$, with larger coordinate balls $B_i\Subset X\setminus Z$. Apply the local assertion to $\mu_i=\mathbf1_{K_i}\mu$ on $B_i$. After a positive affine change of its bounded local potential, we obtain $v_i\in\PSH(B_i)$ with
\[
 1\leq v_i\leq2,\qquad
 \mu_i\ll(\ddc v_i)^n.
\]
Zero pieces may be omitted. Pick a smooth $\chi_i$ on $B_i$ which is zero near $K_i$ and $-3$ near $\partial B_i$. For a sufficiently small $\epsilon_i>0$,
\[
 \theta+\ddc\bigl(\rho+\epsilon_i(v_i+\chi_i)\bigr)
 \geq\tfrac\kappa2\omega+\epsilon_i\ddc v_i.
\]
Inside $B_i$ take the maximum of this potential with $\rho$, and extend it by $\rho$ outside $B_i$. The two pieces glue because $v_i+\chi_i\leq-1$ on a boundary collar. Finally take the maximum with $-C_i$, where $C_i$ is large enough that this last maximum does not change the function near $K_i$. Constants are $\theta$-plurisubharmonic since $\theta\geq0$. We have obtained a global bounded $q_i\in\PSH(X,\theta)$ which agrees with $\rho+\epsilon_i v_i$ near $K_i$. Positivity of bounded mixed products therefore gives
\[
 \mathrm{MA}_\theta(q_i)\geq\epsilon_i^n(\ddc v_i)^n
 \quad\hbox{near }K_i,
 \qquad \mu_i\ll\mathrm{MA}_\theta(q_i).
\]

Translate each $q_i$ so that $q_i\leq0$, and choose positive numbers $t_i$ with
\[
 \sum_i t_i<1,\qquad
 \sum_i t_i\|q_i\|_\infty<\infty.
\]
The uniformly convergent sum $\psi=\sum_i t_iq_i$ is bounded and $\theta$-plurisubharmonic. For a fixed index $i$ it is a convex combination
$\psi=t_iq_i+(1-t_i)w_i$ with a bounded $w_i\in\PSH(X,\theta)$: use the remaining terms and the unused weight on the constant zero potential. Hence
\[
 \mathrm{MA}_\theta(\psi)\geq t_i^n\mathrm{MA}_\theta(q_i).
\]
Every $\mu_i$ is absolutely continuous with respect to this single finite measure. The sets $K_i$ cover $X\setminus Z$, and $\mu(Z)=0$. Thus $\mu\ll\mathrm{MA}_\theta(\psi)$. The Radon--Nikodym theorem supplies the asserted integrable density. If $\mu=0$, simply take $\psi=0$ and $f=0$.
\end{proof}

\Needspace{15\baselineskip}
\begin{theorem}\label{thm:general-measure-minimal}
Let $\theta\geq0$ be smooth and big, and assume $v_-(\theta)>0$. Let $\mu$ be a finite non-pluripolar measure with $\mu(X)>0$.
For every $\lambda>0$ there is a unique $u\in\mathcal E_0$ satisfying
\[
 \mathrm{MA}_\theta^{\mathrm{ext}}(u)=e^{\lambda u}\mu.
\]
There also exist $v\in\mathcal E_0$, normalized by $\sup_Xv=0$, and a unique positive scalar $c$ such that
\[
 \mathrm{MA}_\theta^{\mathrm{ext}}(v)=c\mu.
\]
The second assertion does not claim uniqueness of the normalized potential. No bounded mass property is assumed.
\end{theorem}
\begin{proof}
The reference $0$ is model because $\theta\geq0$, and $P_{\epsilon\theta}(0)=0$ for every $\epsilon>0$. Every bounded $\theta$-plurisubharmonic function belongs to $\mathcal E_0$, since a sufficiently negative constant is a competitor for each of its stretched envelopes. The lower estimate of \cref{thm:relative-uniform-mass-lower}, with $\epsilon\downarrow0$, gives
\begin{equation}\label{eq:minimal-lower-mass}
 \mathrm{MA}^{\mathrm{ext}}_\theta(w)(X)\geq v_-(\theta)
 \qquad(w\in\mathcal E_0).
\end{equation}

\emph{Positive exponent.} Write $\mu=f\nu$ by the preceding lemma, with $\nu=\mathrm{MA}_\theta(\psi)$ and $\psi$ bounded. Let $f_j=\min(f,j)$ for positive integers $j$. These densities are bounded and nonzero as measures. The bounded subsolution theorem gives bounded $u_j$ with
\[
 \mathrm{MA}_\theta(u_j)=e^{\lambda u_j}f_j\nu.
\]
Exponential order shows that $u_j$ decreases, because $f_j$ increases. In particular $u_j\leq u_1$. On the other hand
\[
 v_-(\theta)\leq\mathrm{MA}_\theta(u_j)(X)
       \leq e^{\lambda\sup_Xu_j}\mu(X).
\]
Thus their suprema have a common lower bound. The decreasing limit $u$ is a genuine $\theta$-plurisubharmonic function, and convergence holds in $L^1(X)$. Moreover
\[
 \mathrm{MA}_\theta(u_j)\leq e^{\lambda\sup_Xu_1}\mu.
\]
Apply \cref{thm:relative-fixed-measure-continuity} with reference zero. It gives $u\in\mathcal E_0$ and weak convergence of the Monge--Amp\`ere measures to $\mathrm{MA}^{\mathrm{ext}}_\theta(u)$. The right sides converge in total variation to $e^{\lambda u}f\nu$: the functions $e^{\lambda u_j}f_j$ converge almost everywhere for $\nu$ and are bounded by $e^{\lambda\sup_Xu_1}f\in L^1(\nu)$. This proves the equation. Its measure is finite, so \cref{cor:relative-finite-exponential-order} gives uniqueness.

\Needspace{10\baselineskip}
\emph{Zero exponent.} For each positive integer $j$, let $u_j\in\mathcal E_0$ be the just-constructed solution with exponent $1/j$. We first bound the normalization. Let $h$ be the bounded solution of
\[
 \mathrm{MA}_\theta(h)=e^h\min(f,1)\nu,
\]
supplied by the bounded subsolution theorem. Set $H=\|h\|_\infty$ and $w_j=h+(j-1)H$. Then $w_j/j\geq h$, and
\[
 e^{-w_j/j}\mathrm{MA}_\theta(w_j)
 \leq e^{-h}\mathrm{MA}_\theta(h)
 \leq\mu
 =e^{-u_j/j}\mathrm{MA}^{\mathrm{ext}}_\theta(u_j).
\]
Exponential order, now with exponent $1/j$, gives $u_j\leq w_j$. Put
\[
 M_j=\sup_Xu_j,\qquad v_j=u_j-M_j,
 \qquad c_j=e^{M_j/j}.
\]
Then $\sup_Xv_j=0$ and $0<c_j\leq e^H$. Equation \eqref{eq:minimal-lower-mass} also gives
\[
 v_-(\theta)\leq c_j\int_X e^{v_j/j}\,\mathrm{d}\mu
                 \leq c_j\mu(X),
\]
so $c_j\geq v_-(\theta)/\mu(X)>0$. After extraction, compactness of normalized potentials gives $v_j\to v$ in $L^1(X)$, with $v$ genuine and normalized, and $c_j\to c>0$.

By \cref{lem:normalized-exponential-npp}, $a_j=e^{v_j/j}\to1$ in $L^1(\mu)$. Consequently
\[
 \mathrm{MA}^{\mathrm{ext}}_\theta(v_j)=c_ja_j\mu
       \longrightarrow c\mu
 \quad\hbox{in total variation},\qquad
 \mathrm{MA}^{\mathrm{ext}}_\theta(v_j)\leq e^H\mu.
\]
Apply \cref{thm:relative-fixed-measure-continuity} again. It gives $v\in\mathcal E_0$ and identifies the limit as $\mathrm{MA}^{\mathrm{ext}}_\theta(v)$. The scalar assertion in \cref{cor:relative-finite-exponential-order}, applied in both directions to two solutions, proves uniqueness of $c$.
\end{proof}

Both equations in the minimal class are obtained using domination by one fixed finite measure. This result also provides comparison potentials for a prescribed reference: finite-measure exponential order can be applied with reference zero and any second potential more singular than zero. The proof in \cref{sec:general-measure-relative} uses this observation after its first approximation with moving bounded reference measures.


\Needspace{18\baselineskip}
\subsection{The maximal coefficient under positive lower mass}

For a Hermitian background $h$, the scalar in the preceding equation equals the subsolution threshold $b_h(f)$ from \cref{sec:individual-density-subsolutions}. One envelope-class solution bounds the coefficient of every subsolution by scalar comparison. Dominated continuity also identifies the limits of the normalized exponential solutions.

\begin{corollary}\label{cor:positive-lower-mass-exact-threshold}
Let $h$ be Hermitian with $v_-(h)>0$, and let $0\leq f\in L^1(\mathrm{d}V)$ be nonzero. There is $v\in\mathcal E_{\mathrm{env}}(X,h,0)$ such that
\[
 \mathrm{MA}^{\mathrm{ext}}_h(v)=b_h(f)f\,\mathrm{d}V.
\]
The number $b_h(f)>0$ is the unique positive scalar for which this equation has a solution in that envelope class.

Let $u_\lambda$ be the greatest exponential solution from \cref{thm:individual-exponential-subsolution}, and set $v_\lambda=u_\lambda-\sup_Xu_\lambda$. Each $u_\lambda$ belongs to $\mathcal E_{\mathrm{env}}(X,h,0)$. If $\lambda_j\downarrow0$ and $v_{\lambda_j}\to v$ in $L^1$, then the convergence holds in capacity, $v$ belongs to the same envelope class and satisfies the displayed equation.
\end{corollary}
\begin{proof}
By \cref{thm:general-measure-minimal}, there are $c>0$ and $w\in\mathcal E_{\mathrm{env}}(X,h,0)$ with $\mathrm{MA}^{\mathrm{ext}}_h(w)=cf\,\mathrm{d}V$. If $s\in\PSH(X,h)$ satisfies $\mathrm{MA}^{\mathrm{ext}}_h(s)\geq af\,\mathrm{d}V$ with $a>0$, then
\[
 \mathrm{MA}^{\mathrm{ext}}_h(w)
       \leq\frac ca\,\mathrm{MA}^{\mathrm{ext}}_h(s).
\]
Every $h$-plurisubharmonic $s$ satisfies $s\preceq0$. Thus the scalar assertion of \cref{cor:relative-extended-exponential-order}, with reference zero, gives $c/a\geq1$. Maximizing over subsolution coefficients yields $b_h(f)\leq c$. The reverse inequality holds because $w$ is itself a subsolution. The same argument identifies the scalar of every solution in this envelope class.

For $\lambda>0$, \cref{thm:general-measure-minimal} supplies an exponential solution $s_\lambda$ in the envelope class. The greatest-solution property gives $s_\lambda\leq u_\lambda$. Upward stability, \cref{prop:relative-elementary}, places $u_\lambda$ and its constant translate $v_\lambda$ in that class.

Write $c_\lambda=\exp(\lambda\sup_Xu_\lambda)$. By \cref{cor:individual-positive-threshold-exponential}, $c_\lambda\to b_h(f)>0$ and
\[
 \mathrm{MA}^{\mathrm{ext}}_h(v_\lambda)
       =c_\lambda\mathrm{e}^{\lambda v_\lambda}f\,\mathrm{d}V
       \longrightarrow b_h(f)f\,\mathrm{d}V
       \quad\text{in total variation}.
\]
Since $v_\lambda\leq0$ and $c_\lambda$ stays bounded for small $\lambda$, these measures are dominated by one finite non-pluripolar measure $C f\,\mathrm{d}V$. Apply \cref{thm:relative-fixed-measure-continuity} to the stated $L^1$ convergent sequence. It gives class membership, capacity convergence and the displayed equation for its limit.
\end{proof}

\Cref{thm:normalized-compactness} supplies such convergent subsequences. The final section gives examples with more than one normalized zero-exponent solution; accordingly, the corollary identifies cluster equations and does not assert convergence of the whole family.


\section{Equations with a model reference}
\label{sec:general-measure-relative}

We now prove \cref{thm:als-general-measure}. Write the given measure as $\mu=f\nu$, where $\nu$ is the Monge--Amp\`ere measure of a bounded potential $\psi$. To keep the reference $\varphi$, we first smooth $\psi$: the prescribed ordinary-type theorem applies to the approximating data, exponential comparison provides a common lower envelope, and moving-measure continuity preserves the equation in the limit. With $\psi$ fixed, we then approximate the integrable density $f$ by bounded densities. Here the solutions for the reference zero give upper bounds, while positive lower mass prevents collapse.

\begin{proof}[of \cref{thm:als-general-measure}]
The proof has four steps: continuous positive densities, bounded densities, integrable densities, and finally the zero exponent.

Let $\mu$ be the given nonzero finite non-pluripolar measure. Apply \cref{lem:bounded-measure-carrier} to the Hermitian, hence semipositive and big, background $\omega$ to write
\[
 \mu=f\nu,\qquad \nu=\omega_\psi^n,\qquad
 \psi\in\PSH(X,\omega)\cap L^\infty(X),\qquad
 0\leq f\in L^1(\nu).
\]
\emph{Step 1: positive continuous density.} Fix $\lambda>0$ and suppose first that $f$ is continuous and strictly positive. By \cref{thm:bounded-exact-regularization}, choose smooth strictly $\omega$-plurisubharmonic $\psi_j\downarrow\psi$. These potentials are uniformly bounded; write $\|\psi_j\|_\infty\leq L$ and $\nu_j=\omega_{\psi_j}^n$. Bounded Bedford--Taylor continuity gives $\nu_j\rightharpoonup\nu$. The measure $f\nu_j$ has a continuous positive density with respect to smooth volume, so it belongs to every finite $L^p$ space. Hence the prescribed ordinary-type theorem \cref{thm:als-prescribed-lp} gives $u_j\sim\varphi$ with
\[
 \mu_{u_j}=\mathrm{e}^{\lambda u_j}f\nu_j.
\]

We first bound the $u_j$ from above, uniformly in $j$. Choose $K$ with
$\mathrm{e}^{-\lambda(\psi_j+K)}\leq\min_Xf$ for all $j$. Then
\[
 \mathrm{e}^{-\lambda(\psi_j+K)}\omega_{\psi_j+K}^n
 \leq f\nu_j
 =\mathrm{e}^{-\lambda u_j}\mu_{u_j}
 \leq \mathrm{e}^{-\lambda u_j}\mathrm{MA}^{\mathrm{ext}}_\omega(u_j),
\]
where the last inequality follows from $\omega\geq\theta$ and positivity of extended mixed products. Apply \cref{cor:singular-exponential-order} in the background $\omega$, with first potential $\psi_j+K$ and second potential $u_j$. Its hypotheses hold: the first potential is bounded, so it admits a smaller-background envelope, and $u_j\preceq\psi_j+K$. Thus $u_j\leq\psi_j+K\leq L+K=:C$.

By \eqref{eq:relative-normalized-model-bound}, $u_j\leq\varphi+C$. Since
\[
 \mu_{u_j}\leq\|f\|_\infty \mathrm{e}^{\lambda u_j}\omega_{\psi_j}^n,
\]
\cref{prop:exponential-common-rooftop} gives a single $h\in\mathcal E_\varphi$ below all $u_j$. The bounds $h\leq u_j\leq\varphi+C$ and \cref{thm:normalized-compactness} give, after extraction, an $L^1$ limit $u\not\equiv-\infty$. The moving-measure theorem \cref{thm:relative-moving-measure-continuity} now applies, because $\psi_j\to\psi$ in capacity and $\mu_{u_j}\leq \mathrm{e}^{\lambda C}\|f\|_\infty\nu_j$. It gives $u\in\mathcal E_\varphi$, convergence in capacity and $\mu_{u_j}\rightharpoonup\mu_u$.

The right-hand sides converge weakly to $\mathrm{e}^{\lambda u}f\nu$. Indeed, $\mathrm{e}^{\lambda u_j}-\mathrm{e}^{\lambda u}\to0$ in capacity with a common bound, so its integral against $\nu_j$ tends to zero by the uniform capacity estimate for these measures. The fixed function $\mathrm{e}^{\lambda u}$ is bounded and quasi-continuous. Approximating it by a continuous function outside an open set of arbitrarily small capacity, bounded product continuity and the same capacity estimate allow this fixed weight in $\nu_j\rightharpoonup\nu$. Multiplication by the continuous function $f$ and by a continuous test function preserves both estimates. Hence $\mu_u=\mathrm{e}^{\lambda u}f\nu$.

\emph{Step 2: bounded density.} Suppose next that $0\leq f\in L^\infty(\nu)$ and $\int_Xf\,\mathrm{d}\nu>0$. Choose positive continuous functions $f_j$, uniformly bounded by a constant $F$, with $f_j\to f$ almost everywhere for $\nu$. Such a sequence is obtained from the density of continuous functions in $L^1(\nu)$ by clipping to a common interval, adding positive constants tending to zero, and extracting a summably convergent subsequence.

There is a bounded nonnegative function $g$ with $g\nu\ne0$ and $g\leq f_j$ for all $j$. Indeed, choose $a>0$ with $\nu(\{f\geq2a\})>0$. By Egorov's theorem, there is a measurable subset $E\subseteq\{f\geq2a\}$ of positive measure on which the convergence is uniform; hence $f_j\geq a$ on $E$ for all $j\geq N$, for some $N$. The finitely many remaining functions $f_1,\ldots,f_{N-1}$ are positive and continuous, so they have positive minima on $X$. If $b>0$ is a lower bound for these minima and for $a$, then $g=b\mathbf1_E$ works.

By Step 1, let $u_j\in\mathcal E_\varphi$ solve $\mu_{u_j}=\mathrm{e}^{\lambda u_j}f_j\nu$. The theorem for the reference zero, \cref{thm:general-measure-minimal}, applied to the finite nonzero measure $g\nu$, gives $s\in\mathcal E_0$ satisfying
\[
 \mathrm{MA}^{\mathrm{ext}}_\theta(s)=\mathrm{e}^{\lambda s}g\nu.
\]
Apply \cref{cor:relative-extended-exponential-order} with reference zero, first potential $s\in\mathcal E_0$ and second potential $u_j\preceq0$. Since $g\nu\leq f_j\nu$, it gives $u_j\leq s$. With $C=\sup_Xs$ we have
\[
 \mu_{u_j}\leq \mathrm{e}^{\lambda C}F\nu,
 \qquad
 \kappa\leq \mathrm{e}^{\lambda\sup_Xu_j}F\nu(X),
\]
where the second inequality uses $\mu_{u_j}(X)\geq\kappa$. These two bounds on the suprema give a subsequence converging in $L^1$ to some $u\not\equiv-\infty$. The fixed-measure continuity theorem \cref{thm:relative-fixed-measure-continuity} puts $u$ in $\mathcal E_\varphi$ and identifies the weak limit of $\mu_{u_j}$ as $\mu_u$.

By \cref{lem:fixed-npp-integrals}, $\mathrm{e}^{\lambda u_j}\to \mathrm{e}^{\lambda u}$ in $L^1(\nu)$, and $f_j\to f$ in $L^1(\nu)$ by dominated convergence. Since both factors are uniformly bounded,
\[
 \mathrm{e}^{\lambda u_j}f_j\longrightarrow \mathrm{e}^{\lambda u}f
 \quad\hbox{in }L^1(\nu).
\]
Thus the right-hand sides converge in total variation, which proves the equation for $f$.

\emph{Step 3: integrable density.} For general $f\in L^1(\nu)$, apply Step 2 to $f_j=\min(f,j)$ and obtain $u_j\in\mathcal E_\varphi$ with $\mu_{u_j}=\mathrm{e}^{\lambda u_j}f_j\nu$. Since $f_j\leq f_{j+1}$, exponential order gives $u_{j+1}\leq u_j$. Furthermore
\[
 \kappa\leq \mathrm{e}^{\lambda\sup_Xu_j}\mu(X),
 \qquad \mu_{u_j}\leq \mathrm{e}^{\lambda\sup_Xu_1}\mu.
\]
Hence the decreasing limit $u$ is not identically $-\infty$, and the fixed-measure continuity theorem gives $u\in\mathcal E_\varphi$ and $\mu_{u_j}\rightharpoonup\mu_u$. Dominated convergence gives $\mathrm{e}^{\lambda u_j}f_j\to \mathrm{e}^{\lambda u}f$ in $L^1(\nu)$, so $\mu_u=\mathrm{e}^{\lambda u}\mu$. Uniqueness follows from \cref{cor:relative-extended-exponential-order}, applied in both directions.

\emph{Step 4: zero exponent.} For each positive integer $j$, Step 3 with exponent $1/j$ gives $u_j\in\mathcal E_\varphi$ with $\mu_{u_j}=\mathrm{e}^{u_j/j}\mu$. The zero-exponent part of the proof of \cref{thm:general-measure-minimal} applies to these $u_j$, with the same comparison potential $s\in\mathcal E_0$, the solution of $\mathrm{MA}^{\mathrm{ext}}_\theta(s)=\mathrm{e}^s\mu$. Indeed, the comparison $u_j\leq s+(j-1)\max(0,\sup_Xs)$ made there uses exponential order with reference zero, first potential in $\mathcal E_0$ and second potential $u_j\preceq0$; the lower mass bound $v_-(\theta)$ is replaced by $\kappa$; and \cref{thm:relative-fixed-measure-continuity} is applied with reference $\varphi$. We obtain $v\in\mathcal E_\varphi$ with $\sup_Xv=0$ and $c>0$ with $\mu_v=c\mu$. Finally, the scalar assertion of \cref{cor:relative-extended-exponential-order}, applied in both directions, shows that $c$ is unique.
\end{proof}


\subsection{Nonuniqueness in the minimal envelope class}

In this subsection the manifold and the background are no longer fixed. The examples of \cref{uni-thm:analytic-example} are smooth solutions on projective manifolds with nonclosed Hermitian backgrounds. Comparison with a K\"ahler form shows that these backgrounds have both positive lower mass and bounded mass, and smooth potentials belong to the minimal envelope class. Thus even both mass bounds together allow several normalized solutions with the same scalar.

\begin{corollary}\label{cor:analytic-nonuniqueness-envelope-class}
In every complex dimension $n\geq3$, there are a projective manifold $X$ of dimension $n$, a smooth nonclosed Hermitian form $\omega$ on $X$, and a nonzero finite non-pluripolar measure $\nu$ on $X$ such that
\[
 0<v_-(\omega)\leq v_+(\omega)<+\infty,
\]
but the equation
\[
 \mathrm{MA}^{\mathrm{ext}}_\omega(u)=c\nu,
 \qquad u\in\mathcal E_{\mathrm{env}}(X,\omega,0),
 \qquad \sup_Xu=0,
\]
has at least two smooth solutions $u$ for $c=1$, and no solution for any other positive constant $c$. Relative to a K\"ahler volume form on $X$, the measure $\nu$ has a real analytic density, positive off two connected smooth real hypersurfaces and vanishing quadratically along them.
\end{corollary}
\begin{proof}
Use the form and the two potentials of \cref{uni-thm:analytic-example}. Since the manifold is projective, \cref{lem:kahler-modification-bounded-mass} gives bounded mass. Fix a K\"ahler form $\Omega$ on $X$ and choose $\varepsilon>0$ with $\varepsilon\Omega\leq\omega$. Lower-volume monotonicity, \cref{thm:volume-monotonicity}, applied to the zero potential in both backgrounds, gives
\[
 v_-(\omega)\geq v_-(\varepsilon\Omega)
 =\varepsilon^n\int_X\Omega^n>0.
\]
Here the potentials of the same ordinary type as zero are exactly the bounded ones, so these lower volumes are the infima defining $v_-$.

The two smooth potentials belong to $\mathcal E_{\mathrm{env}}(X,\omega,0)$ by \cref{prop:relative-elementary}, and their extended measures are the ordinary top powers. Their common measure $\nu$ is finite and does not charge pluripolar sets. The construction gives $c=1$, and \cref{cor:relative-extended-exponential-order} shows that no other positive scalar is possible. The density and its zero set were computed in \cref{uni-thm:analytic-example}.
\end{proof}

